Orbiter contains a knowledge base of mathematical data. 
This is like a database, except for the way that the data is stored. 
The data does not reside in files, like in a database. 
Rather, the data is compiled into the Orbiter executable. 
This means that the data is immediately available, without any delay. 
We will now discuss what data exists in Orbiter and how to access it.

\bigskip



In most cases, 
the data that Orbiter stores includes the automorphism group of the objects.
Individual objects can be accessed by creating them from the catalogue.
Other functions exist that can write reports about all objects of a certain type.


\bigskip


The command 

\sourcecode{make_table_of_surfaces}

creates a table of all cubic surfaces contained in the Orbiter knowledge base.
The output is shown in Appendix~\ref{sec:data:cubic:surfaces}.


\bigskip

The command 

\sourcecode{make_table_of_quartic_curves}

creates a table of all quartic curves contained in the Orbiter knowledge base.
%The output is shown in Appendix~\ref{sec:data:cubic:surfaces}.


\bigskip

It is possible to export data on cubic surfaces and quartic curves. 
Two file types are supported: csv and sql. 
For instance, 
the command 

\sourcecode{cubic_surfaces_table_q17}

exports the cubic surfaces over the field $\bbF_{17}$ 
contained in the Orbiter knowledge base.
The files 
\verb'table_of_cubic_surfaces_q17_info.csv' and 
\verb'table_of_cubic_surfaces_q17_data.sql' are created.
For each isomorphism class of cubic surfaces, one line is generate in the file.
The command 

\sourcecode{quartic_curves_table_q17}

exports all quartic curves over the field $\bbF_{17}$ 


\bigskip

